\documentclass[12pt]{article}
\usepackage[utf8]{inputenc}
\usepackage[T1]{fontenc}
\usepackage[margin=1in]{geometry}
\usepackage{setspace}
\usepackage{xcolor}
\usepackage[colorlinks=true, urlcolor=blue!60!black, linkcolor=blue!60!black]{hyperref}
\setlength{\parindent}{0em}
\setlength{\parskip}{0.45em}
\renewcommand{\baselinestretch}{1.0}
\begin{document}
\begin{center}
{\Large \textit{National Pork Producers Council v. Ross}, 598 U.S. 356 (2023)\par}
\vspace{0.3em}
{\normalsize Case Brief --- One-Page Edition\par}
\end{center}

\textbf{Memory Jogger.} Proposition 12 and pork: no ``almost per se'' extraterritoriality rule (\href{https://www.law.cornell.edu/supremecourt/text/21-468}{598 U.S.\ 356}).

\textbf{Facts.} California's Proposition 12 barred in-state sales of whole pork from breeding pigs ``confined in a cruel manner.'' Pork-industry groups alleged that California's market share would force producers nationwide to restructure, while conceding the law does not discriminate.

\textbf{Procedural History.} The district court dismissed; the Ninth Circuit affirmed.

\textbf{Votes.} Affirmed, 5--4 (Gorsuch, with shifting pluralities; Roberts, Alito, Kavanaugh, and Jackson dissenting in part).

\textbf{Holding.} Antidiscrimination lies at the ``very core'' of the Dormant Commerce Clause; there is no ``almost per se'' rule against nondiscriminatory laws with out-of-state effects. \textit{Pike} survives but mainly targets hidden protectionism and burdens on interstate transportation.

\textbf{Reasoning.} \textit{Baldwin} and \textit{Healy} were protectionist price-control cases, and petitioners' rule would imperil ordinary state laws. Three Justices found costs and moral benefits incommensurable; four found no substantial burden pleaded. Roberts would have remanded for \textit{Pike} balancing; Kavanaugh warned of ``California knows best'' balkanization.

\textbf{Examples.} \emph{Same-Side:} a foie gras sales ban. \emph{Opposite-Side:} a law tying in-state prices to out-of-state prices. \emph{Fence-Sitter:} sales bans tied to out-of-state labor policies.

\textbf{Critique.} Originalists praised restraint; critics fear balkanization by large States and an unsettled \textit{Pike}.

\textbf{Key Quotation.} ``Companies that choose to sell products in various States must normally comply with the laws of those various States. Assuredly, under this Court's dormant Commerce Clause decisions, no State may use its laws to discriminate purposefully against out-of-state economic interests. But the pork producers do not suggest that California's law offends this principle. Instead, they invite us to fashion two new and more aggressive constitutional restrictions on the ability of States to regulate goods sold within their borders. We decline that invitation. While the Constitution addresses many weighty issues, the type of pork chops California merchants may sell is not on that list.'' --- Justice Gorsuch, 598 U.S.\ at 364.

\end{document}
